%% kapa: exam=1 round=32 year=2021 period=2 form=A subject=accounting subject_ko=회계학 seq=3 number=43
%% source: Q-Net 감정평가사 기출문제 (cst00309) artlSeq=5208397 file=extracted/1008/5208397/A형/2021년도 제32회 감정평가사 1차 2교시 A형.pdf
%% conversion: kordoc 4.24.51 -> tools/gen1.py (auto); review=proofread-sonnet; flags=none
\begin{question}{43}
공정가치 측정에 관한 설명으로 옳지 않은 것은?
\begin{choices}
  \item 공정가치란 측정일에 시장참여자 사이의 정상거래에서 자산을 매도할 때 받거나 부채를 이전할 때 지급하게 될 가격이다.
  \item 공정가치는 시장에 근거한 측정치이며 기업 특유의 측정치가 아니다.
  \item 공정가치를 측정하기 위해 사용하는 가치평가기법은 관측할 수 있는 투입변수를 최소한으로 사용하고 관측할 수 없는 투입변수를 최대한으로 사용한다.
  \item 기업은 시장참여자가 경제적으로 최선의 행동을 한다는 가정 하에, 시장참여자가 자산이나 부채의 가격을 결정할 때 사용할 가정에 근거하여 자산이나 부채의 공정가치를 측정하여야 한다.
  \item 비금융자산의 공정가치를 측정할 때는 자신이 그 자산을 최고 최선으로 사용하거나 최고 최선으로 사용할 다른 시장참여자에게 그 자산을 매도함으로써 경제적 효익을 창출할 수 있는 시장참여자의 능력을 고려한다.
\end{choices}
\end{question}
\answer{3}
